% Generated by Naberius. Content is escaped; shell escape is disabled.
\documentclass[12pt,a4paper]{letter}
\usepackage{fontspec}
\defaultfontfeatures{Ligatures=TeX}
\setmainfont[Path=fonts/,BoldFont=lmroman10-bold.otf,ItalicFont=lmroman10-italic.otf,BoldItalicFont=lmroman10-bolditalic.otf]{lmroman10-regular.otf}
\usepackage[spanish]{babel}
\usepackage[margin=25.0mm,headheight=16pt]{geometry}
\usepackage{graphicx,amsmath,amssymb,tabularx,booktabs}
\usepackage{fancyhdr}
\usepackage[unicode,hidelinks]{hyperref}
\hypersetup{pdftitle={El oficio de escribir},pdfauthor={Documento de demostración},pdfkeywords={}}
\pagestyle{fancy}
\fancyhf{}\fancyhead[L]{\small }\fancyfoot[C]{\small \thepage}\renewcommand{\headrulewidth}{0pt}
\fancypagestyle{plain}{\fancyhf{}\fancyhead[L]{\small }\fancyfoot[C]{\small \thepage}\renewcommand{\headrulewidth}{0pt}}
\linespread{1.15}
\setlength{\parindent}{0.0mm}
\setlength{\parskip}{7.0pt}
\signature{Documento de demostración}
\address{}
\date{}
\begin{document}
\begin{letter}{Comité académico}
\opening{Estimado comité:}
{\large\bfseries Resumen\par}
Prueba tipográfica de \LaTeX{}: español, áéíóú y referencias.
\fontsize{12.0}{14.399999999999999}\selectfont
{\LARGE\bfseries El oficio de escribir\par}

{\large Del primer borrador a una idea que se entiende.\par}

\par\medskip\phantomsection\label{b3C9EB2A7-151E-452C-898C-C2E310F8CC02}{\large\bfseries Una idea necesita espacio\par}

Escribir es dar forma a lo que pensamos. Este documento de demostración muestra cómo organizar un texto, enlazar sus fuentes y cuidar su presentación sin interrumpir la escritura.

\par\medskip\phantomsection\label{b4444FE3B-8BAB-49E9-B630-0BED1B6EE8B5}{\large\bfseries Estructura antes que apariencia\par}

Usa los estilos de la barra para expresar qué significa cada párrafo. El esquema se construye a partir de tus títulos; la plantilla se ocupa de la tipografía.

Inserta una tabla, una figura o una ecuación desde el menú Insertar.

\begin{itemize}
\item Guarda tu trabajo como .nab: contenido, imágenes y referencias viajan juntos.

\end{itemize}
\par\medskip\phantomsection\label{b748A2195-4C7F-43CE-83F6-7100733DFB03}{\large\bfseries Un documento, dos maneras de verlo\par}

Alterna entre páginas y lienzo continuo. El texto y el historial son los mismos; la paginación definitiva corresponde al PDF compuesto con \LaTeX{}.

\LaTeX{} y \TeX{} conservan su composición tipográfica.

\par\medskip\phantomsection\label{bF75E1FE4-C928-4AEC-B19B-06EED4C4FCEC}{\large\bfseries Datos complementarios\par}

\par\medskip\noindent\rule{\linewidth}{0.4pt}\par\medskip

\par\medskip\phantomsection\label{n1347B508-3672-47ED-9C14-5B63C1697D76}{\small Tabla 1. Resultados de prueba\par}
\begin{tabularx}{\linewidth}{|>{\raggedright\arraybackslash}X|>{\raggedright\arraybackslash}X|}\hline
\textbf{Variable} & \textbf{Valor} \\ \hline
 &  \\ \hline
\end{tabularx}\par{\small }

\begin{equation}\sqrt{\alpha  + \beta }\label{n6A3148D6-BC0C-4681-BC51-6FE3E1F2A357}\end{equation}

\closing{Atentamente,}
\end{letter}
\end{document}